\documentclass[12pt]{article}

% Preambulo
\usepackage[frenchb]{babel}
\usepackage{float}
\usepackage[utf8]{inputenc}  
\usepackage{lmodern}
\usepackage[T1]{fontenc}  
\usepackage{graphicx}
\usepackage{indentfirst}
\usepackage{courier}
\usepackage{geometry}
\usepackage{glossaries}
\usepackage{eurosym}
\newcommand {\cod} [1] { \texttt {#1}} 
\newcommand {\p} [1] { \paragraph {#1}}
\usepackage{hyperref}
\hypersetup{
    colorlinks,
    citecolor=black,
    filecolor=black,
    linkcolor=black,
    urlcolor=black
}
\usepackage[urldate=comp,dateabbrev=false,sorting=none]{biblatex}
\bibliography{bibliography}

\title{Rapport de Stage 2A}
\author{Antônio Guilherme \textsc{Ferreira Viggiano}}
%\date{}
\makeindex

\geometry{paperwidth=210mm,paperheight=297mm,
top=30mm,bottom=30mm,
left=30mm,right=30mm}
% Documento

\makeglossaries

\begin {document}


\begin{titlepage}
\begin{center}
% Upper part of the page

\begin{figure}[ht]
  \begin{minipage}[b]{0.5\linewidth}
    \centering
    \includegraphics[width=0.7\textwidth]{img/logo_ecm}
    \label{logo_ecm}
  \end{minipage}
  \hspace{0.5cm}
  \begin{minipage}[t]{0.5\linewidth}
    \vspace{-4cm}
    \begin{center}
    \includegraphics[width=0.4\textwidth]{img/logo_asert}
    \end{center}
    \label{logo_asert}
  \end{minipage}
\end{figure}

%\textsc{\LARGE Intelligence Artificielle et Prolog}\\[1.5cm]


% Title
\vspace{2cm}
\hrule  ~\\[0.4cm]
{ \huge \bfseries Stage assistant ingénieur en développement logiciel}\\[0.4cm]

\hrule  ~\\[0.5cm]

\textsc{\LARGE Rapport de stage deuxième année}\\[1.5cm]


\end{center}


\vfill


\large\emph{Elève :} \hfill Antônio Guilherme \textsc{Ferreira Viggiano} ~\\

\large\emph{Tuteur pédagogique :} \hfill M. Pascal \textsc{Préa}

\large\emph{Tuteur de stage :} \hfill M. Clemenceau Roberto \textsc{Da Silva}

\vfill

\begin{center}
  \large Stage realisé du 3 Juin au 26 Juillet \\
  (8 semaines)
\end{center}






\end{titlepage}


%Glossary
\renewcommand*{\glspostdescription}{. Page:}
\renewcommand*{\glossaryname}{}

\pagebreak	
\tableofcontents
%!TEX root = index.tex

%Entradas do glossário
\newglossaryentry{EV}
{
  name=EV,
  description={-- Electric Vehicle. A vehicle that uses one or more electric motors or traction motors for propulsion. Electric vehicles include cars, trains, trucks, planes, boats, motorcycles, among others. The term is frequently interchanged with electric cars, designating battery-only cars or hybrid electric cars},
  plural=EVs
}

\newglossaryentry{energy-mix}
{
  name=energy mix,
  description={-- The many energy sources of a region, either renewable or non-renewable. Theses may include coal, biomass, diesel, wave power, wind energy, etc}
}

\newglossaryentry{well-to-wheel}
{
  name=well to wheel,
  description={-- The life-cycle analysis used to evaluate environmental impacts associated with all stages of a vehicle's life. This would include in particular gas emissions in the process of powering up batteries and transporting them}
}

\newglossaryentry{plug-in-hybrid}
{
  name=plug-in hybrid,
  description={-- Hybrid vehicle that uses batteries which can be fully recharged by connecting a plug to an external power source, either at home or at a power station. This is used to reduce costs related to battery replacement and to enable all-electric operation, reducing greenhouse gas emissions},
  plural=plug-in hybrids
}

\newglossaryentry{Better-Place}
{
  name=Better Place,
  description={-- An international company that developed and sold battery-switching battery for electric cars. The corporation did not succeed mainly because of non-standardization of battery models, and has declared bankruptcy in Israel and has initiated the proceedings in Denmark},
}


\newglossaryentry{frunk}
{
  name=frunk,
  description={-- The front trunk. Electric vehicles are usually able to have both rear and front trunks because the batteries are located in the floor, and the motor between the drive wheels},
}

\newglossaryentry{hybrid}
{
	name=hybrid car,
	description={-- Vehicle that uses more than one power source to move the wheels. Most commonly, it refers to hybrid electric vehicles, which combine internal combustion engines and electric motors}
}

\newglossaryentry{li-ion}
{
	name=lithium-ion battery,
	description={-- A rechargeable battery used to power electric vehicles. Li-ion batteries have high energy density and slow self-discharge when not in use, and are known for being more expensive than older nickel metal batteries},
	plural=lithium-ion batteries
}

\newglossaryentry{range-extender}
{
	name=range extender,
	description={-- A system used in hybrid cars where the internal combustion engine is used solely to charge up the battery, and not to power the wheels. This system is used as an emergency tool, when battery is running out of charge},
	plural=range extenders
}

\newglossaryentry{range-anxiety}
{
	name=range anxiety,
	description={-- An expression used to designate a driver's worry that his vehicle won't reach its destination on the current charge, especially because of battery charge restrictions and limited electric power stations on the way}
}

\newglossaryentry{zero-emission}
{
	name=zero-emission vehicle,
	description={-- A vehicle that has zero gas emissions at the tailpipe. Emissions from the manufacturing process are whatsoever ignored in this definition},
	plural=zero tailpipe emissions
}

\newglossaryentry{gas-emission}
{
	name=greenhouse gas,
	description={-- Gas substance (such as water vapor or carbon dioxide) that absorb part of the infrared radiation emitted by the surface of Earth. This process decreases heat loss, keeping the planet warm. Human activities such as industrial or tailpipe emissions increase greenhouse gases concentration on air, changing Earth's average temperature},
	plural=greenhouse gas emissions
}

\newglossaryentry{fuel-cell}
{
	name=fuel cell,
	description={-- Device capable of converting chemical energy into electricity. An electrolyte initiates a chemical reaction that generates electric power when fuel and oxidants, such as hydrogen and oxygen, flow into it. While no greenhouse gases are produced in the process, hydrogen production and distribution may have a significant environmental impact}
}

\newglossaryentry{government-incentives}
{
	name=government incentives,
	description={-- Subsidies or government policies beneficial for electric cars, which vary with federal government and vehicle model. Most incentives are tax credits and personal income tax deduction for costumers},
	plural=government policy
}

\newglossaryentry{bev}
{
	name=battery electric vehicle,
	description={-- Also referred to as all-electric vehicle, it is a vehicle that is use a battery as the only mean to store electricity needed to power the motor},
	plural=battery electric vehicles
}



\pagebreak
\section{Remerciements}

\section*{Remerciements}

Je voudrais remercier à tous ceux qui ont contribué pour le bon déroulement de ce stage, en particulier à :

\vspace{1em}

\begin{description}
\item[Mme Mirta JUAREZ et Mme Cécile LOUBET --] pour les réponses rapides à toute question concernant le stage à l'international.
\item[M. Marco Aurélio RABELO --] président-directeur général de Asert, pour être si proche des employés.
\item[M. Clemenceau Roberto DA SILVA --] directeur de technologie de l'entreprise, qui m'a bien accueilli par email et personnellement au sein de l'entreprise.
\item[M. Thiago RODRIGUES PEREIRA --] architect logiciel, pour m'avoir formé dans le framework de l'entreprise et m'avoir introduit aux bonnes pratiques en génie logiciel.
\end{description}


\pagebreak
\section{Glossaire}
\printglossary

\pagebreak
%!TEX root = index.tex
\printbibliography[heading=bibintoc]


\end{document}
